% Adapted from Fall 2022 Final Q4b (Advent of Code)

\begin{blocksection}
\question Implement \texttt{find\_pair}, which takes a list of integers \texttt{s} and returns the two elements of \texttt{s} that sum to 2020, in the order that they appear in \texttt{s}. Assume that exactly one pair of elements sums to 2020. You may call \lstinline{list.contains(s, x)}, which returns whether \texttt{s} contains \texttt{x}.
\begin{lstlisting}
pub fn find_pair(s: List(Int)) -> List(Int) {
  case s {
    [] -> []
    [x, ..rest] ->
      case ______________________________ {
        True -> ________________
        False -> ________________
      }
  }
}
assert find_pair([979, 1721, 366, 299, 675, 1456]) == [1721, 299]
assert find_pair([2000, 20]) == [2000, 20]
assert find_pair([1, 1010, 3, 1010]) == [1010, 1010]
\end{lstlisting}
\begin{solution}
\begin{lstlisting}
pub fn find_pair(s: List(Int)) -> List(Int) {
  case s {
    [] -> []
    [x, ..rest] ->
      case list.contains(rest, 2020 - x) {
        True -> [x, 2020 - x]
        False -> find_pair(rest)
      }
  }
}
\end{lstlisting}
\end{solution}
\begin{questionmeta}
  % Instead of trying every pair, we only ask one question about each element \texttt{x}: does its partner \lstinline{2020 - x} appear later in the list? If so, we already know both numbers, so we can return \lstinline{[x, 2020 - x]} without finding where the partner is.
  % We only search \texttt{rest} rather than all of \texttt{s}. Any partner earlier in the list would have been found when we were at that earlier element, so by the time we reach \texttt{x}, nothing before it can be part of the answer. Searching \texttt{rest} also means \texttt{x} can never be paired with itself, which is why the last assert works even though 1010 + 1010 is 2020. (The original exam guaranteed that \texttt{s} contains no 1010, since its Python version could pair an element with itself.)
  % The \lstinline{[] -> []} case is never reached when a pair exists, but Gleam requires every case to be handled.
  % This takes quadratic time: each call to \lstinline{list.contains} can scan the rest of the list, and there can be one call per element.
\end{questionmeta}
\end{blocksection}